% =====================================================================
%  PETITION FOR WRIT OF MANDAMUS (MINISTERIAL) AND COMPLAINT FOR
%  DECLARATORY JUDGMENT
%  Circuit Court for Baltimore County, Maryland
%  McLaughlin & Fernandez v. Bank of New York Mellon, et al.
%
%  PRO SE DRAFT.  Bracketed placeholders [ ... ] must be completed.
%  NOT FOR FILING WITHOUT REVIEW BY A LICENSED MARYLAND ATTORNEY.
% =====================================================================
\documentclass[12pt,letterpaper]{article}

\usepackage[margin=1in]{geometry}
\usepackage{mathptmx}
\usepackage{setspace}
\usepackage{enumitem}
\usepackage{fancyhdr}
\usepackage{booktabs}
\usepackage{longtable}
\usepackage{tabularx}
\usepackage{array}
\usepackage{parskip}
\usepackage{microtype}
\usepackage{xcolor}
\usepackage[hidelinks]{hyperref}
\usepackage{ragged2e}
\usepackage{lastpage}

% ---------------------------------------------------------------- layout
\renewcommand{\thesection}{\Roman{section}}
\renewcommand{\thesubsection}{\Alph{subsection}}
\renewcommand{\thesubsubsection}{\arabic{subsubsection}}
\setcounter{secnumdepth}{3}
\setcounter{tocdepth}{2}

\pagestyle{fancy}
\fancyhf{}
\fancyhead[L]{\footnotesize McLaughlin \& Fernandez v. Bank of New York Mellon, et al.}
\fancyhead[R]{\footnotesize Pet. for Writ of Mandamus}
\fancyfoot[C]{\footnotesize Page \thepage\ of \pageref{LastPage}}
\renewcommand{\headrulewidth}{0.4pt}

% ------------------------------------------------- paragraph numbering
\newcounter{para}
\newcommand{\np}{%
  \par\medskip
  \refstepcounter{para}%
  \noindent\hangindent=0.45in\hangafter=1
  \textbf{\arabic{para}.}\hspace{0.15in}\ignorespaces}

% ------------------------------------------------- shorthand
\newcommand{\CJP}{Md.\ Code Ann., Cts.\ \&\ Jud.\ Proc.}
\newcommand{\ET}{Md.\ Code Ann., Est.\ \&\ Trusts}
\newcommand{\R}{\S}

\begin{document}

% =====================================================================
%  DRAFT NOTICE -- DELETE BEFORE FILING
% =====================================================================
\thispagestyle{empty}
\begin{center}
\fbox{\parbox{0.92\textwidth}{\centering\textbf{DRAFT --- NOT FOR FILING\\[4pt]
\small This is an unverified drafting instrument. Every bracketed item in the
form [ ... ] must be completed and every citation verified against the current
Maryland Rules and Code before any filing. This document was prepared with
automated drafting assistance and is not legal advice.}}}
\end{center}
\vspace{0.35in}

% =====================================================================
%  CAPTION
% =====================================================================
\begin{center}
\textbf{\large IN THE CIRCUIT COURT FOR BALTIMORE COUNTY, MARYLAND}
\end{center}
\vspace{2pt}\hrule\vspace{8pt}

\noindent
\begin{minipage}[t]{0.54\textwidth}
\raggedright
CAUSTIN LEE McLAUGHLIN\\
1000 North Woodington Road\\
Baltimore, Maryland 21229\\[5pt]
\hspace*{0.25in}\textit{and}\\[5pt]
JUAN JOSE FERNANDEZ\\
{[Street address]}\\
{[City, State, ZIP]}\\[8pt]
\hspace*{0.25in}Petitioners,\\[8pt]
\hspace*{0.25in}v.\\[8pt]
THE BANK OF NEW YORK MELLON;\\
DELAWARE TRUST COMPANY;\\
CENTRUS ENERGY CORP.\ (f/k/a USEC INC.);\\
COMPUTERSHARE TRUST COMPANY, N.A.\\[8pt]
\hspace*{0.25in}Respondents.
\end{minipage}\hfill
\begin{minipage}[t]{0.42\textwidth}
\raggedright
Case No.: \underline{\hspace{1.45in}}\\
(New Action)\\[12pt]
\textbf{PETITION FOR WRIT OF MANDAMUS (MINISTERIAL) AND COMPLAINT FOR DECLARATORY JUDGMENT}\\[10pt]
Filed Pro Se\\[4pt]
Request for Waiver of Prepayment of Costs\\
filed herewith (Md.\ Rule 1-325)
\end{minipage}

\vspace{8pt}\hrule
\vspace{10pt}

% =====================================================================
%  TABLE OF CONTENTS
% =====================================================================
{\singlespacing\tableofcontents}
\newpage

% =====================================================================
%  PRELIMINARY STATEMENT
% =====================================================================
\section*{Preliminary Statement}
\addcontentsline{toc}{section}{Preliminary Statement}

\np This is a petition for a writ of mandamus of the ministerial variety and,
in the alternative, a complaint for declaratory judgment. Petitioners Caustin
Lee McLaughlin and Juan Jose Fernandez ask this Court to compel Respondents to
perform acts that Petitioners contend are non-discretionary, ministerial, and
required by [the governing trust instruments / the applicable transfer-agent and
indenture provisions], and to declare the respective rights and obligations of
the parties.

\np Petitioners expressly do not ask this Court to direct any Respondent to
exercise, or to refrain from exercising, discretion committed to that Respondent
by law or by contract. The acts for which relief is sought are framed as
non-discretionary acts of record-keeping, acknowledgment, processing,
accounting, and transmittal of completed appointment packages.

\np Petitioners further ask this Court to compel the processing and final
consideration of Schedule A and Schedule C appointment packages for both
Petitioners, as one component of a larger community-infrastructure and
community-resilience implementation described in Part IV.C and Exhibit A.

% =====================================================================
\section{Nature of the Action}
% =====================================================================

\np This action arises from (i) the alleged failure of Respondents, in their
capacities as [trustee / successor trustee / indenture trustee / transfer agent
--- specify per Respondent], to perform ministerial acts in respect of
Petitioner McLaughlin's record holdings and [trust / beneficial] interests, and
(ii) the alleged failure of the appropriate appointing authority to process
completed Schedule A and Schedule C appointment packages for Petitioners.

\np Petitioners contend that the duties at issue are purely ministerial: they
are duties that leave nothing to judgment once the facts are established, and
that are owed to Petitioners as a matter of record. Petitioners contend that
they have a clear legal right to performance, that Respondents owe a clear,
non-discretionary duty of performance, and that no adequate remedy at law
exists. See Part V, \textit{infra}.

\np Petitioners additionally seek declaratory relief under the Maryland Uniform
Declaratory Judgments Act, \CJP\ \R\R\ 3-401 \textit{et seq.}, to determine the
parties' rights and obligations, including the legal sufficiency of the
tendered appointment packages and the scope of the ministerial duties at issue.

% =====================================================================
\section{The Parties}
% =====================================================================

\subsection{Petitioners}

\np \textbf{Caustin Lee McLaughlin} is an individual residing at 1000 North
Woodington Road, Baltimore, Maryland 21229. Mr.\ McLaughlin proceeds \textit{pro
se}. He is, on information and belief, a debtor-in-possession in an active
Chapter 11, Subchapter V bankruptcy proceeding, Case No.\ [24-2383], and he
asserts a medically and genomically documented neurodevelopmental disability
that requires cognitive accommodation in court proceedings. The supporting
clinical and genomic records are attached as Exhibit [\,], and are referenced
solely to establish status for the purposes of \CJP\ \R\ [\,] and the federal
excepted-service appointment authority discussed in Part VI.F.

\np \textbf{Juan Jose Fernandez} is an individual residing at [Street address,
City, State, ZIP] and is the spouse of Mr.\ McLaughlin. Mr.\ Fernandez seeks
appointment to a legal-technical support position as described in Exhibit A.

\np No attorney has entered an appearance for either Petitioner. Petitioners
request that all filings be served upon them at the addresses stated above.

\subsection{Respondents}

\np \textbf{The Bank of New York Mellon} is a New York banking corporation with
its principal office at [principal office address --- verify]. It is named as a
Respondent in its capacity as [trustee / successor trustee / indenture trustee
--- verify and specify the instrument by date and parties].

\np \textbf{Delaware Trust Company} is a Delaware trust company with its
principal office at [principal office address --- verify]. It is named as a
Respondent in its capacity as [successor trustee / trustee --- verify].

\np \textbf{Centrus Energy Corp.}, formerly known as USEC Inc., is a Delaware
corporation with its principal office at [principal office address --- verify].
It is named as a Respondent in its capacity as [issuer / obligor / successor by
reorganization --- verify].

\np \textbf{Computershare Trust Company, N.A.} is a federally chartered trust
company with its principal office at [principal office address --- verify]. It
is named as a Respondent in its capacity as [transfer agent / registrar /
successor transfer agent --- verify].

\np The precise capacity of each Respondent, the identity of the appointing
authority for the Schedule A and Schedule C positions at issue, and the
instrument or agreement giving rise to the alleged ministerial duties are
matters of record that Petitioners will establish by the exhibits listed at the
end of this Petition. Where a capacity is stated above in brackets, that
capacity must be completed before filing.

% =====================================================================
\section{Jurisdiction and Venue}
% =====================================================================

\np \textbf{Subject-matter jurisdiction.} The Circuit Court for Baltimore
County has jurisdiction over this action under \CJP\ \R\ 1-501, which confers
upon the circuit courts of this State full common-law and equity jurisdiction,
and under \CJP\ \R\ 1-502. Jurisdiction to issue a writ of mandamus in a
non-administrative matter is governed by Maryland Rule 15-701. Administrative
mandamus, which is not sought here, is governed by Maryland Rules 7-401 through
7-403.

\np \textbf{Declaratory jurisdiction.} Jurisdiction to grant declaratory relief
is conferred by \CJP\ \R\R\ 3-401 \textit{et seq.}, and in particular \R\R\
3-406 and 3-409(c). A court may grant declaratory relief with or without
coercive relief, and may grant declaratory relief whether or not a writ of
mandamus issues. \CJP\ \R\ 3-409(c).

\np \textbf{Venue.} Venue is proper in Baltimore County under \CJP\ \R\
6-201(a), because [Respondent(s)] carr[y/ies] on a regular business in
Baltimore County, and under \CJP\ \R\ 6-201(b), because there is more than one
Respondent, at least one of whom could be sued in Baltimore County, and because
the cause of action arose in Baltimore County. In the alternative, venue is
proper under \CJP\ \R\ 6-201(b) because a substantial part of the events giving
rise to the claims occurred in Baltimore County at Petitioner McLaughlin's
residence and place of business.

\np \textbf{Prepayment of costs.} Petitioners are unable to pay the costs of
this proceeding and intend to file a Request for Waiver of Prepayment of Costs
under Md.\ Rule 1-325. See also \CJP\ \R\ 7-201. Petitioners request that the
clerk accept this Petition for filing upon submission of the Request, and that
the Court rule upon the Request before requiring any payment.

% =====================================================================
\section{Statement of Facts}
% =====================================================================

\subsection{The Instruments and the Alleged Trust Chain}

\np Petitioner McLaughlin is, on information and belief, the holder and/or
beneficiary of record of interests recorded in [name of the instrument or
account], administered by one or more of Respondents. The instrument(s) at
issue are: [describe each indenture, trust agreement, transfer-agent agreement,
or account agreement by date, parties, and governing law].

\np Under those instruments, the trustee, successor trustee, indenture trustee,
and/or transfer agent named therein owes duties of record-keeping,
acknowledgment, registration of transfer, accounting, and transmittal. Those
duties are, Petitioners contend, ministerial as to the specific acts identified
in Part VI.A.

\np On information and belief, the relevant corporate relationships are as
follows: [state, for each Respondent, the executed instrument and the
succession history by date]. Respondent Centrus Energy Corp.\ is the successor
by reorganization to USEC Inc.\ [verify].

\subsection{Petitioner McLaughlin's Proceedings and Status}

\np Mr.\ McLaughlin has proceeded \textit{pro se} in [number] matters in this
State and in the federal courts, including [list case numbers and courts]. Five
consolidated petitions for a writ of certiorari and a related petition for a
writ of mandamus and prohibition were filed in the Supreme Court of the United
States on or about September 6, 2026. A related proceeding is pending under
Case No.\ [24-C-26-003970] in the Circuit Court for Baltimore City.

\np Mr.\ McLaughlin has been granted leave to proceed \textit{in forma
pauperis} in [identify courts], and that status is a matter of record. He
asserts a documented neurodevelopmental disability affecting executive function
and auditory processing under conditions of stress, and requests written
instructions and reasonable accommodation in these proceedings. See Exhibit
[\,]; 42 U.S.C.\ \R\ 12131 \textit{et seq.}; 28 C.F.R.\ \R\ 35.130(b)(7).

\subsection{The Proposed Introduction of Petitioners Under a Larger Implementation}

\np Petitioners have proposed, and seek to implement, an integrated
community-infrastructure and community-resilience program. That program is
described in the materials attached as Exhibit [\,] and is referred to in this
Petition as the ``Integrated Implementation.'' The Integrated Implementation
includes, among other components, [brief neutral description of the components
that Petitioners intend to pursue, e.g., a community resilience and housing
component, a records-and-trust administration component, and a personnel
component].

\np A necessary personnel component of the Integrated Implementation is the
placement of two qualified individuals --- Petitioners McLaughlin and Fernandez
--- into technical and legal-technical support roles. Petitioners contend that
they are qualified for the positions described in Exhibit A and that the
appointment packages for those positions are complete.

\np On or about October 21, 2026, a hiring and onboarding schedule was prepared
proposing the retention of both Petitioners in specialized technical and legal
support roles: Mr.\ McLaughlin as a Systems Help Desk Lead / Technical
Specialist, and Mr.\ Fernandez as a Jurisprudential Technician / Legal IT
Coordinator. See Exhibit A.

\np Petitioners do not contend that they are entitled to appointment as a matter
of right. They contend only that the processing, acknowledgment, transmittal, and
final consideration of completed appointment packages is a non-discretionary
act, and that the failure to perform that act is remediable by mandamus.

\subsection{The Failure to Act Giving Rise to This Petition}

\np On or about [date], Petitioners submitted to Respondents a written request
for [describe the ministerial act demanded, e.g., acknowledgment of record,
registration of transfer, an account of the trust, transmittal of a completed
appointment package]. A copy is attached as Exhibit [\,].

\np Respondents have failed to perform the requested act, and have failed to
respond within [number] days, notwithstanding [number] written follow-ups dated
[dates]. See Exhibits [\,]--[\,].

\np On or about [date], Petitioners submitted the completed Schedule A and
Schedule C appointment packages to the appropriate appointing authority. The
packages have not been processed to final consideration. See Exhibit A.

\np No Respondent has asserted, to Petitioners' knowledge, any ground of
substantive objection to the requests; the failure is, Petitioners contend, a
failure to perform a ministerial act of record, not a reasoned exercise of
discretion.

% =====================================================================
\section{The Governing Legal Standard for a Ministerial Writ of Mandamus}
% =====================================================================

\np A writ of mandamus is an extraordinary remedy, and it is not issued as a
matter of course. Maryland Rule 15-701. It is available to compel the
performance of a ministerial duty, but it is not available to control or direct
the exercise of discretion committed to the respondent.

\np To obtain a writ of mandamus, a petitioner must show (1) that the petitioner
has a clear legal right to the relief sought; (2) that the respondent owes a
clear legal duty to perform the act demanded; and (3) that the petitioner has no
other adequate remedy at law. \textit{Brack v.\ Wells}, 184 Md.\ 86, 90--91, 40
A.2d 319, 321 (1944); \textit{Philip Morris Inc.\ v.\ Angeletti}, 358 Md.\ 689,
712, 752 A.2d 200, 212 (2000).

\np A duty is ministerial when it is a duty the respondent is required to
perform upon a given state of facts, in a prescribed manner, in obedience to a
mandate of legal authority, without regard to the respondent's own judgment
upon the propriety of the act to be done. Conversely, where the respondent is
vested with discretion to decide whether, and in what manner, to act, mandamus
does not lie. \textit{Falls Rd.\ Cmty.\ Ass'n v.\ Baltimore County}, 439 Md.\ 1
(2014) (declining to issue a writ of mandamus where the challenged decision was
committed to official discretion).

\np A circuit court may order an agency or official to take an action that the
officer's office requires the officer to take, but may not order an action that
is discretionary. The writ may also be used to correct a public record. See Md.\
Courts, \textit{Mandamus} (describing conventional and administrative mandamus).

\np Because Petitioners seek relief of the conventional (non-administrative)
kind, Maryland Rules 7-401 through 7-403, which govern administrative mandamus
and judicial review of quasi-judicial agency orders, do not supply the
procedural framework here.

% =====================================================================
\section{Argument}
% =====================================================================

\subsection{The Acts Sought Are Ministerial, Not Discretionary}

\np The acts for which Petitioners seek relief are acts of record:
acknowledgment of a document, entry of a record, registration of a transfer,
production of an account, and transmittal of a completed package. Each of these
acts is, Petitioners contend, required to be performed upon presentation of a
completed and conforming request, in a prescribed manner, without regard to the
Respondent's own judgment upon the propriety of the act. That is the definition
of a ministerial duty.

\np Petitioners do not ask this Court to direct any Respondent to approve,
deny, or grant anything. Petitioners ask only that the act of \textit{processing
and responding} --- an act from which the exercise of judgment is separate --- be
compelled.

\subsection{Petitioners Have a Clear Legal Right}

\np Petitioners' right to performance is, they contend, clear and not dependent
on any contested issue of fact reserved for adjudication. The right arises from
[the instruments] and from [the applicable appointment regulations]. See Exhibit
[\,].

\subsection{Petitioners Have No Adequate Remedy at Law}

\np An ordinary action at law for damages would not produce the specific record
act that Petitioners require, and would not cure the structural impediment to
the Integrated Implementation. Declaratory relief alone, without coercion, would
not secure performance. Petitioners therefore contend that no other adequate
remedy exists.

\subsection{Count I --- Mandamus to Compel Ministerial Acts}

\np Petitioners incorporate the foregoing paragraphs and seek a writ of mandamus
compelling each Respondent, to the extent it is the proper party, to perform the
specific ministerial acts identified in Exhibit [\,], including: (i) to
acknowledge receipt of Petitioners' conforming request; (ii) to enter and
maintain the corresponding record; (iii) to produce an account of the interests
at issue; and (iv) to transmit the completed package to the appropriate
appointing authority.

\subsection{Count II --- Declaratory Judgment}

\np Petitioners seek a declaration under \CJP\ \R\R\ 3-401 \textit{et seq.}
determining: (i) the identity of the proper party for each ministerial act;
(ii) whether the tendered requests and packages are legally sufficient and
conforming; and (iii) the respective rights and obligations of the parties with
respect to the acts described in Count I.

\subsection{Count III --- Mandamus to Compel Processing of the Schedule A and Schedule C Packages}

\np Petitioners further seek a writ of mandamus compelling the completion of
final consideration of the Schedule A and Schedule C appointment packages
described in Exhibit A.

\np \textbf{The appointment authorities.} Schedule A of the excepted service,
5 C.F.R.\ \R\ 213.3102, including the authority for the appointment of persons
with disabilities at 5 C.F.R.\ \R\ 213.3102(u), permits appointment without
competitive examination. Schedule C, 5 C.F.R.\ \R\R\ 213.3301--213.3302,
governs positions of a confidential or policy-determining character. Both are
excepted-service appointing authorities. Petitioners contend that Mr.\
McLaughlin is eligible for consideration under 5 C.F.R.\ \R\ 213.3102(u), and
that the positions described in Exhibit A may be filled under Schedule A or
Schedule C as the appointing authority directs. See also 5 U.S.C.\ \R\R\ 2301,
2302 (merit system principles and prohibited personnel practices).

\np \textbf{The limited scope of the relief sought.} Petitioners acknowledge
that appointment is committed to the discretion of the appointing authority and
that no provision of law confers upon them a right to be appointed. Petitioners
seek only the ministerial act of processing, acknowledging, and bringing the
completed packages to final consideration within a time certain.

\np \textbf{Reservation.} To the extent this Court concludes that it lacks
authority to compel any part of the processing described in this Count because
the appointing authority is not before the Court or because the act is
committed to federal discretion, Petitioners request that this Court so state in
a written order, so that the question may be preserved for review and, if
appropriate, transferred or refiled in the correct forum.

\subsection{The Integrated Implementation}

\np The personnel component described in Count III is one component of the
Integrated Implementation described in Part IV.C. Petitioners seek a
declaration under Count II that the personnel component is not an unlawful
component of the Integrated Implementation, and that the ministerial acts
required to advance it may proceed.

% =====================================================================
\section{Request for an Order to Show Cause and a Hearing}
% =====================================================================

\np Petitioners request that the Court issue an order requiring each Respondent
to show cause, within a time fixed by the Court, why the writ should not issue,
and that the Court set this matter for a hearing. A proposed form of order to
show cause is attached as Exhibit B. Petitioners request that the hearing be
conducted with the accommodations described in Part IV.B.

% =====================================================================
\section{Prayer for Relief}
% =====================================================================

WHEREFORE, Petitioners respectfully request that this Court:

\begin{enumerate}[label=(\alph*),leftmargin=0.6in,itemsep=3pt]
\item issue a writ of mandamus compelling each Respondent, to the extent it is
the proper party, to perform the ministerial acts identified in Count I and
Exhibit [\,];
\item issue a writ of mandamus compelling the processing and final
consideration of the Schedule A and Schedule C appointment packages described in
Count III and Exhibit A;
\item enter a declaratory judgment under \CJP\ \R\R\ 3-401 \textit{et seq.}
determining the parties' rights and obligations as requested in Count II;
\item grant a waiver of the prepayment of costs under Md.\ Rule 1-325;
\item set this matter for a hearing on an order to show cause; and
\item grant such other and further relief as the Court deems just and proper.
\end{enumerate}

\vspace{0.25in}
\noindent Respectfully submitted,

\vspace{0.35in}
\noindent\rule{2.8in}{0.4pt}\\
Caustin Lee McLaughlin\\
1000 North Woodington Road\\
Baltimore, Maryland 21229\\
Petitioner, Pro Se

\vspace{0.35in}
\noindent\rule{2.8in}{0.4pt}\\
Juan Jose Fernandez\\
{[Street address]}\\
{[City, State, ZIP]}\\
Petitioner, Pro Se

\vspace{0.2in}
\noindent Dated: \underline{\hspace{0.4in}} day of
\underline{\hspace{1.4in}}, 2026.

% =====================================================================
\section{Verification}
% =====================================================================

I, \underline{\hspace{2.6in}}, being first duly sworn, do hereby verify that
the facts stated in the foregoing Petition for Writ of Mandamus and Complaint
for Declaratory Judgment are true and correct to the best of my knowledge,
information, and belief, except as to matters stated upon information and
belief, which I believe to be true.

\vspace{0.3in}
\noindent\rule{2.8in}{0.4pt}\\
Caustin Lee McLaughlin

\vspace{0.3in}
\noindent\rule{2.8in}{0.4pt}\\
Juan Jose Fernandez

\vspace{0.3in}
\noindent
\begin{tabular}{@{}p{3.2in}p{2.6in}@{}}
State of Maryland, County of Baltimore & \\[6pt]
Subscribed and sworn to before me this & \\[6pt]
\underline{\hspace{0.6in}} day of \underline{\hspace{1.2in}}, 2026. & \\[6pt]
& \rule{2.5in}{0.4pt}\\
& Notary Public\\
& My commission expires: \underline{\hspace{1.2in}}\\
\end{tabular}

% =====================================================================
\section{Certificate of Service}
% =====================================================================

I hereby certify that on this \underline{\hspace{0.4in}} day of
\underline{\hspace{1.4in}}, 2026, a copy of the foregoing Petition for Writ of
Mandamus (Ministerial) and Complaint for Declaratory Judgment was served upon
each Respondent by [certified mail, return receipt requested / personal service
/ the Court's electronic filing system] at the address of its registered agent
for service of process in this State, as follows:

\vspace{0.15in}
\noindent
\begin{tabularx}{\textwidth}{@{}p{2.4in}X@{}}
\toprule
Respondent & Service address / registered agent\\
\midrule
The Bank of New York Mellon & [address --- verify]\\
Delaware Trust Company & [address --- verify]\\
Centrus Energy Corp. & [address --- verify]\\
Computershare Trust Company, N.A. & [address --- verify]\\
\bottomrule
\end{tabularx}

\vspace{0.3in}
\noindent\rule{2.8in}{0.4pt}\\
Caustin Lee McLaughlin, Petitioner, Pro Se

% =====================================================================
\section*{Exhibit Index}
\addcontentsline{toc}{section}{Exhibit Index}
% =====================================================================

\begin{longtable}{@{}p{0.9in}p{5.3in}@{}}
\toprule
Exhibit & Description\\
\midrule
\endhead
A & Schedule A / Schedule C Hiring Implementation Matrix for Petitioners
McLaughlin and Fernandez (see page after next).\\
B & Proposed Order to Show Cause (see final page).\\
C & [Clinical and genomic records supporting Petitioner McLaughlin's asserted
disability status.]\\
D & [The instrument(s), indenture(s), trust agreement(s), or account
agreement(s) at issue.]\\
E & [Petitioners' written request for the ministerial act, with proof of
submission.]\\
F & [Respondents' responses or notices of non-response.]\\
G & [Description of the Integrated Implementation.]\\
\bottomrule
\end{longtable}

\newpage
% =====================================================================
\section*{Exhibit A --- Schedule A / Schedule C Hiring Implementation Matrix}
\addcontentsline{toc}{section}{Exhibit A --- Hiring Implementation Matrix}
% =====================================================================

\np This Exhibit states the positions for which Petitioners seek processing and
final consideration. The appointing authority and the specific series and grade
must be completed before filing.

\vspace{0.15in}
\noindent
\begin{longtable}{@{}p{1.1in}p{1.5in}p{1.75in}p{1.3in}@{}}
\toprule
Candidate & Position title & Appointing authority & Series / grade\\
\midrule
\endhead
McLaughlin & Systems Help Desk Lead / Technical Specialist & Sch.\ A, 5 C.F.R.\
\R\ 213.3102(u) & [Series/Grade]\\
Fernandez & Jurisprudential Technician / Legal IT Coordinator & Sch.\ A, 5
C.F.R.\ \R\ 213.3102, or Sch.\ C, 5 C.F.R.\ \R\ 213.3301 & [Series/Grade]\\
\bottomrule
\end{longtable}

\vspace{0.2in}
\noindent
\begin{longtable}{@{}p{1.5in}p{2.6in}p{1.5in}@{}}
\toprule
Phase & Milestone & Target date\\
\midrule
\endhead
Phase 1 & Offer acceptance and credentialing & Week 1\\
Phase 2 & Background and suitability processing & Weeks 2--3\\
Phase 3 & Systems access and onboarding & Weeks 4--5\\
Phase 4 & Integration into the Integrated Implementation & Weeks 6--8\\
\bottomrule
\end{longtable}

\vspace{0.15in}
\noindent\footnotesize Source: Petitioners' hiring and onboarding schedule
dated October 21, 2026. The duties, series, grade, and appointing authority for
each position must be confirmed by the appointing authority before filing.

\newpage
% =====================================================================
\section*{Exhibit B --- Proposed Order to Show Cause (Proposed Only)}
\addcontentsline{toc}{section}{Exhibit B --- Proposed Order to Show Cause}
% =====================================================================

\begin{center}
\textbf{\textsc{Proposed Order}}\\[4pt]
\textit{(This is a proposed form. It has no effect unless and until signed by a
judge of this Court.)}
\end{center}

\vspace{0.2in}
Upon consideration of the Petition for Writ of Mandamus (Ministerial) and
Complaint for Declaratory Judgment filed by Petitioners Caustin Lee McLaughlin
and Juan Jose Fernandez, it is, this \underline{\hspace{0.4in}} day of
\underline{\hspace{1.2in}}, 2026,

\vspace{0.15in}
\textbf{ORDERED} that Respondents The Bank of New York Mellon; Delaware Trust
Company; Centrus Energy Corp.\ (f/k/a USEC Inc.); and Computershare Trust
Company, N.A., shall show cause, if any they have, on or before
\underline{\hspace{1.2in}}, why the relief requested in the Petition should not
be granted; and it is further

\textbf{ORDERED} that this matter is set for hearing on
\underline{\hspace{1.2in}}, at \underline{\hspace{0.8in}} a.m./p.m., in
Courtroom \underline{\hspace{0.6in}} of the Circuit Court for Baltimore County,
Maryland.

\vspace{0.4in}
\noindent\rule{3.2in}{0.4pt}\\
Judge, Circuit Court for Baltimore County

\end{document}
